\section*{Overview}

Bitmask DP appears when the state is a small subset and the transition depends on which elements have already been
taken. The size limit is strict, but within that range the technique is extremely expressive.

\section*{When to Use It}

Use bitmask DP when:

\begin{itemize}
  \item \(n\) is small enough that \(2^n\) states are plausible,
  \item the state is ``which items are already used'',
  \item transitions depend on the last chosen item, the next free item, or a submask split.
\end{itemize}

\section*{Core Idea}

Represent the chosen set by a mask. Then define a DP state such as:
\[
dp[mask][last] = \text{best answer after visiting exactly the items in } mask \text{ and ending at } last.
\]

\section*{Key Insight}

The technique works because the mask compresses an exponential family of subsets into array indices. That turns
``visited-set'' recursion into an iterative DP over masks.

\section*{Operations / Main Technique}

\begin{itemize}
  \item choose a subset state,
  \item iterate masks in increasing order,
  \item transition by adding one new bit or splitting into submasks.
\end{itemize}

\paragraph{Worked example.}
In TSP-style DP, from state \((mask, last)\), try every city \(\texttt{nxt}\) not in \(\texttt{mask}\) and transition
to \((mask \cup \{nxt\}, nxt)\).

\section*{Correctness Intuition}

Each mask encodes exactly which decisions have already been made. The transition adds one legal next decision, so every
valid configuration is reached from smaller subsets, and every DP value is built from the relevant previous states.

\section*{Complexity Analysis}

Many standard forms run in \(O(n 2^n)\) or \(O(n^2 2^n)\), with \(O(n 2^n)\) memory.

\section*{Implementation}

The code shows an assignment-style bitmask DP where \(\texttt{dp[mask]}\) tracks the best cost after assigning the
first \(\texttt{popcount(mask)}\) rows to the chosen columns.

\section*{Common Pitfalls}

\begin{itemize}
  \item Starting bitmask DP when \(n\) is already too large.
  \item Forgetting whether the transition depends on set size, last item, or both.
  \item Writing \(O(3^n)\) subset loops when \(O(n 2^n)\) would be enough.
\end{itemize}

\section*{Variants / Extensions}

\begin{itemize}
  \item TSP / Hamiltonian path DP.
  \item SOS DP and subset convolution.
  \item State-compression DP on grids.
\end{itemize}

\section*{Practice Problems}

\begin{itemize}
  \item Assignment or matching on \(n \le 20\).
  \item Traveling salesman on a small graph.
  \item Partition a small set into compatible groups.
\end{itemize}

\section*{References}

\begin{itemize}
  \item \href{https://usaco.guide/gold/dp-bitmasks}{USACO Guide: Bitmask DP}.
  \item \href{https://codeforces.com/catalog}{Codeforces Catalog}.
\end{itemize}
